\section{Downstream Benchmark}
\label{sec:experiments}

We evaluate our model, \Ours{} (a grammar-conditioned ESMC$\rightarrow$LLaDA
diffusion language model), on a unified panel of immune-receptor and
protein--protein interaction downstream tasks. The row labelled \OursCkpt{} is
the retrained integrated 7-layer checkpoint; its features are extracted
\emph{post-LLaDA} following the mandatory protocol in \Cref{sec:interface}: the
full record is rendered, passed through a clean (noise-free)
ESMC$\rightarrow$LLaDA forward pass, and mean-pooled once over all residue
positions into a single vector. Baseline evidence is labelled throughout as
paper-reported [P], official-artifact re-scored [A], or locally evaluated [L].
When the downloaded data or local head differs from the anchor protocol, [P]
is the primary reference and [L] is not ranked against it. The provenance
registry is stored under \texttt{outputs/external/}. Numbers for \OursCkpt{}
come from the retrained integrated
7-layer checkpoint
(\texttt{grammar\_v2\_esmc300m\_integrated\_llada\_7l}, step 16{,}000),
reported honestly, including cases where \Ours{} trails specialised baselines.

%--------------------------------------------------------------------
\subsection{Model plug-in interface}
\label{sec:interface}

Every task exposes a single string-spec entry point so that any embedder ---
ours or a baseline --- can be swapped in without touching task code. The factory
\texttt{build\_embedder(spec)} routes on the prefix:

\begin{center}
\small
\begin{adjustbox}{max width=\textwidth}
\begin{tabular}{@{}ll@{}}
\toprule
Spec prefix & Meaning \\
\midrule
\texttt{bioseq-llada:/abs/best.pt} & our model, post-LLaDA features (headline) \\
\texttt{grammar:decoder:global:/abs/best.pt} & post-LLaDA whole-record global mean-pool \\
\texttt{grammar:encoder:/abs/best.pt} & ESMC-encoder-only ablation (control) \\
\texttt{esm2\_150m}, \texttt{protbert}, \texttt{sceptr}, \dots & reference / baseline embedders \\
\bottomrule
\end{tabular}
\end{adjustbox}
\end{center}

The same interface is reused by the TCR harness (\texttt{--embedder}), the MINT
PPI harness (\texttt{--model grammar:.../best.pt}), and the antibody/nanobody
generative scripts. Plugging a newly trained checkpoint into the full benchmark
therefore requires only pointing the spec at the new \texttt{best.pt}.

%--------------------------------------------------------------------
\subsection{TCR tasks}
\label{sec:tcr}

\paragraph{T1: TCR--epitope binding (Nature Methods 2025 split).}
Overall (pooled) AUPRC on the official seen/unseen epitope splits
\citep{nathan2025tcrbench}, CDR3$\beta$
track, AS negatives (\Cref{tab:t1}). Baseline rows are transcribed from
Supplementary Table~4 of the anchor paper [P]. They show strong degradation on
\emph{unseen} epitopes. Training-free controls and \OursCkpt{} are local
evaluations [L], separated from the paper rows and not substituted for them.
Additional paper rows are in \Cref{sec:t1-extra}.

\begin{table}[t]
\centering
\caption{T1 TCR--epitope binding on the Nature Methods 2025 benchmark
(CDR3$\beta$, AS negatives, overall AUPRC / AUROC). [P] denotes paper values
and [L] a local diagnostic. Best [P] value per column in \best{bold}. Further
paper variants are in
\Cref{tab:t1-tier1}.}
\label{tab:t1}
\small
\begin{tabular}{@{}lcccc@{}}
\toprule
Method & Seen AUPRC & Unseen AUPRC & Seen AUROC & Unseen AUROC \\
\midrule
ATM-TCR [P]               & \best{0.6957} & \best{0.5208} & \best{0.6480} & \best{0.5218} \\
TEIM [P]                  & 0.6756 & 0.5159 & 0.6025 & 0.5161 \\
NetTCR-2.2 [P]            & 0.6457 & 0.5062 & 0.5880 & 0.5005 \\
ERGO-LSTM (VDJdb) [P]     & 0.6279 & 0.5171 & 0.5771 & 0.5116 \\
epiTCR [P]                & 0.6023 & 0.4879 & 0.5566 & 0.4887 \\
PanPep [P]                & 0.4979 & 0.4771 & 0.5065 & 0.4825 \\
\midrule
CDR3-kNN [L]              & 0.553 & 0.544 & 0.507 & 0.504 \\
Random [L]                & 0.512 & 0.503 & 0.517 & 0.520 \\
\midrule
\OursCkpt{} [L]           & 0.528 & 0.525 & 0.500 & 0.522 \\
\bottomrule
\end{tabular}
\end{table}

\paragraph{T2: TCR clustering.}
We use two complementary bases. \emph{Basis A} is the NAR Genomics \&
Bioinformatics 2025 paired-chain benchmark \citep{nargab2025tcrcluster}.
\Cref{tab:t2-nargab} reports the paper's own Purity/Retention values [P]. A
separate local audit re-scored the released assignments and matched 6/9 methods
within $\pm0.03$; it also exposed and fixed a reversed length-difference lookup
for LD. Because the remaining discrepancies are unresolved, local re-scores are
not the primary baseline. Basis A is baseline-only (its tools do not expose an
embedding, so \Ours{} is not directly comparable).
\emph{Basis B} is the embedding$\rightarrow$clustering protocol of TCREmbedding
\citep{feng2025tcrembedding}: $K$-means ARI/NMI/Purity averaged over
$K\in\{10,\dots,100\}$ on the 9{,}033 CDR3$\beta$ / 25-epitope dataset
(\Cref{tab:t2}). This is a local embedding audit [L], not a transcription of
the TCREmbedding leaderboard. We evaluate four training-free handcrafted encodings:
$k$-mer and one-hot amino-acid composition, the BLOSUM62 mean substitution
profile, and the Atchley five-factor descriptor. The result partially supports
TCREmbedding's finding that simple features rival data-driven PLM embeddings on
CDR3$\beta$: composition encodings ($k$-mer, one-hot) match or beat the general
PLMs (ESM2-150M, ProtBERT) while the descriptor encodings (BLOSUM62, Atchley)
trail; only the specialised TCR models (SCEPTR, TCR-BERT) lead clearly, so the
decisive factor here is TCR-specific modelling rather than handcrafted vs.\
learned features. These comparisons are restricted to this local harness.

\begin{table}[t]
\centering
\caption{T2 Basis A: NAR-GAB 2025 paired-chain clustering. Values are
paper-reported [P] and replace discrepant local re-scores as the primary
baseline.}
\label{tab:t2-nargab}
\small
\begin{tabular}{@{}lcc@{}}
\toprule
Method [P] & Retention & Purity \\
\midrule
HD        & 0.25 & 0.66 \\
LD        & 0.25 & 0.66 \\
iSMART    & 0.25 & 0.66 \\
GIANA     & 0.25 & 0.66 \\
TCRMatch  & 0.09 & 0.82 \\
clusTCR   & 0.09 & 0.99 \\
GLIPH2    & 0.22 & 0.80 \\
DeepTCR   & 0.92 & 0.63 \\
TCRdist3  & 0.44 & 0.42 \\
\bottomrule
\end{tabular}
\end{table}

\begin{table}[t]
\centering
\caption{T2 Basis B [L]: local TCR clustering audit ($K$-means mean over $K$).
SCEPTR uses CDR3$\beta$+TRBV (native partial input); all other rows use
CDR3$\beta$ only. It is not used as a paper-aligned baseline table.}
\label{tab:t2}
\small
\begin{tabular}{@{}lccc@{}}
\toprule
Embedding & ARI & NMI & Purity \\
\midrule
SCEPTR (contrastive TCR) & \best{0.033} & \best{0.159} & \best{0.339} \\
TCR-BERT                 & 0.016 & 0.095 & 0.287 \\
$k$-mer(2)               & 0.010 & 0.079 & 0.268 \\
ESM2-150M                & 0.007 & 0.068 & 0.258 \\
one-hot (composition)    & 0.007 & 0.064 & 0.254 \\
ProtBERT                 & 0.006 & 0.061 & 0.248 \\
BLOSUM62 (mean profile)  & 0.004 & 0.056 & 0.240 \\
Atchley (5-factor)       & 0.002 & 0.047 & 0.233 \\
\midrule
\OursCkpt{}              & 0.019 & 0.119 & 0.294 \\
\bottomrule
\end{tabular}
\end{table}

\paragraph{T3: TCR representation (SCEPTR few-shot, deep 6-pMHC).}
Per-pMHC nearest-neighbour AUROC \citep{nagano2024sceptr} (macro over 6 pMHC,
100 seeds) at increasing reference-set sizes $k$ (\Cref{tab:t3}). All rows in
this table are local protocol-compatible evaluations [L]. At $k{=}200$, an
independent audit compares the six reference methods with the paper's
Supplementary Information; it does not turn the local rows into paper values.
The complementary broad 24-epitope track and 24-way linear-probe view are in
\Cref{sec:t3-extra}.

\begin{table}[t]
\centering
\caption{T3 local evaluation [L]: deep 6-pMHC few-shot NN AUROC (100 seeds,
random${=}0.5$). Best value per column in \best{bold}; no cell is a
paper-transcribed value.}
\label{tab:t3}
\small
\begin{tabular}{@{}lccccc@{}}
\toprule
Method & $k{=}1$ & $k{=}5$ & $k{=}20$ & $k{=}100$ & $k{=}200$ \\
\midrule
SCEPTR (contrastive)        & 0.640 & \best{0.702} & \best{0.738} & \best{0.775} & \best{0.790} \\
TCRdist (official)          & \best{0.647} & 0.701 & 0.733 & 0.764 & 0.777 \\
CDR3-Levenshtein-NN         & 0.626 & 0.671 & 0.696 & 0.722 & 0.733 \\
$k$-mer(3)                  & 0.607 & 0.665 & 0.690 & 0.717 & 0.728 \\
ESM2-150M                   & 0.591 & 0.631 & 0.656 & 0.685 & 0.696 \\
ProtBERT                    & 0.585 & 0.624 & 0.648 & 0.681 & 0.694 \\
TCR-BERT                    & 0.579 & 0.631 & 0.656 & 0.681 & 0.692 \\
\midrule
\OursCkpt{}                 & 0.581 & 0.622 & 0.648 & 0.671 & 0.681 \\
\bottomrule
\end{tabular}
\end{table}

\paragraph{T4: TCR generation (TCRT5 protocol \citep{tcrt5_2025}).}
Setting~A is a local unconditional repertoire audit [L]. For Setting~B, the
paper evaluator contains 13 scored pMHC blocks plus one reserved simulation
item. We therefore re-score the released prediction artifacts on the exact
sparse-13 set [A] in \Cref{tab:t4}; the older 14-pMHC local table is only a
control and is no longer mixed with these baselines. Target-rich held20 and
full-length $\alpha/\beta$ results remain separate local controls in
\Cref{tab:t4-held20,tab:t4-setc}.

\begin{table}[t]
\centering
\caption{T4 TCR generation with evidence sources kept separate. Setting~A is
local [L] ($n{=}5000$). Setting~B uses released predictions re-scored on the
paper's sparse-13 target list [A] ($K{=}100$). $\downarrow$ lower is better.}
\label{tab:t4}
\small
\setlength{\tabcolsep}{4pt}
\begin{adjustbox}{max width=\textwidth}
\begin{tabular}{@{}lc@{\hskip 1.2em}lccccc@{}}
\toprule
\multicolumn{2}{c}{Setting A [L]} & \multicolumn{6}{c}{Setting B sparse-13 [A]} \\
\cmidrule(r){1-2}\cmidrule(l){3-8}
Method & JSD$\downarrow$ & Method & $n$ & F1@100 & $d_{\text{edit}}\downarrow$ & seq-rec & Char-BLEU \\
\midrule
soNNia & \best{0.043} & TCRT5          & 13 & 0.0003 & 4.467 & \best{0.602} & \best{0.701} \\
OLGA   & 0.061 & GRATCR              & 13 & 0.0000 & \best{4.397} & 0.595 & 0.677 \\
\OursCkpt{} & 0.295 & ER-Transformer & 13 & 0.0000 & 6.068 & 0.337 & 0.594 \\
\bottomrule
\end{tabular}
\end{adjustbox}
\end{table}

\begin{table}[t]
\centering
\caption{T4 local control [L] on the target-rich held20 split ($K{=}100$).
$^{\dagger}$For \Ours{} this split has ${\sim}48\%$ training-pair leakage, so
its row is a \emph{seen}, optimistic control only. This table is not mixed with
the sparse-13 paper baseline.}
\label{tab:t4-held20}
\small
\begin{tabular}{@{}lcccc@{}}
\toprule
Method & F1@100 & $d_{\text{edit}}\downarrow$ & seq-rec & Char-BLEU \\
\midrule
TcrDesign-G (official) & \best{0.153} & \best{1.47} & \best{0.870} & \best{0.986} \\
TCRT5 (HF beam)        & 0.086 & 1.82 & 0.834 & 0.947 \\
\OursCkpt{}$^{\dagger}$ & 0.000 & 6.36 & 0.454 & 0.411 \\
\bottomrule
\end{tabular}
\end{table}

\begin{table}[t]
\centering
\caption{T4 Setting C local control [L], full-length $\alpha/\beta$ assembly. \Ours{}
(unconditional, $n{=}500$) attains high per-residue AA validity but only a small
fraction of $\beta$ chains yield an ANARCI-extractable canonical CDR3$\beta$ ---
its honest capability boundary is CDR3-level, not full-length de novo. TcrDesign
(epitope-conditioned, 34 pMHC $K{=}100$) is the full-length SOTA. ``--'' = not
applicable to that mode.}
\label{tab:t4-setc}
\small
\setlength{\tabcolsep}{4pt}
\begin{adjustbox}{max width=\textwidth}
\begin{tabular}{@{}llcccc@{}}
\toprule
Method & cond. & valid-AA ($\beta/\alpha$) & CDR3 extract & F1@100 & seq-rec \\
\midrule
\OursCkpt{} (full-length) & uncond. & 0.998 / 0.996 & 0.038 & -- & -- \\
TcrDesign (full-length)   & epitope & -- & -- & 0.061 & 0.772 \\
\bottomrule
\end{tabular}
\end{adjustbox}
\end{table}

%--------------------------------------------------------------------
\subsection{Protein--protein interaction (MINT GeneralPPI)}
\label{sec:mint}

The primary baseline is transcribed from MINT Source Data, Figure~2c--h and
Table~S1 [P] \citep{ullanat2026mint}. It uses the paper's task-specific metrics
and full protocol (\Cref{tab:mint-cls}). Our local port caps training at 3,000,
uses an IRBench subset in place of Gold-standard PPI, and builds PDB-Bind from
a different mirror. Those runs are therefore local diagnostics [L], not a MINT
reproduction, and are omitted from the paper-aligned comparison. In addition,
the former local regression implementation reported Pearson/RMSE in a
fold-specific transformed target space; those conclusions are withdrawn. The
code now inverse-transforms predictions before raw-unit metrics, but the
corrected baseline rerun is pending (\Cref{sec:mint-caveat}).

\begin{table}[t]
\centering
\caption{MINT paper-reported baselines [P], transcribed from Source Data
Figure~2c--h / Table~S1 (mean $\pm$ standard deviation). Metrics follow the
paper exactly. No local or \Ours{} value is mixed into this table.}
\label{tab:mint-cls}
\small
\begin{tabular}{@{}llcc@{}}
\toprule
Task & Paper metric & MINT [P] & ESM2-650M [P] \\
\midrule
Human PPI       & Accuracy & $0.8796\pm0.0069$ & $0.8704\pm0.0146$ \\
Yeast PPI       & Accuracy & $0.6870\pm0.0106$ & $0.6074\pm0.0084$ \\
Gold-standard PPI & AUPRC & $0.6872\pm0.0033$ & $0.6639\pm0.0065$ \\
Mutational PPI  & AUPRC   & $0.6449\pm0.0408$ & $0.6351\pm0.0596$ \\
SKEMPI          & Pearson & $0.2974\pm0.1312$ & $0.2562\pm0.1618$ \\
PDB-Bind        & Pearson & $0.5930\pm0.0424$ & $0.5759\pm0.0283$ \\
\bottomrule
\end{tabular}
\end{table}

%--------------------------------------------------------------------
\subsection{Antibody tasks}
\label{sec:ab}

\paragraph{CDR infilling (SabDab 10-fold).}
The primary baseline is Ophiuchus-Ab Table~2 [P], average amino-acid recovery
(AAR, \%) over ten folds (\Cref{tab:ab-cdr}). AntiBERTy and AbLang2 are marked
by that paper as having possible train/test overlap; Ophiuchus-Ab is its
no-overlap zero-shot row. Our earlier masked-infilling run used a different
input layout and is not ranked here; the fold-path loader has been fixed and an
exact local rerun remains pending.

\begin{table}[t]
\centering
\caption{Antibody CDR infilling, paper-reported average AAR (\%) from
Ophiuchus-Ab Table~2 [P]. $^{\dagger}$Possible training-data overlap as flagged
by the paper.}
\label{tab:ab-cdr}
\small
\begin{tabular}{@{}lccc@{}}
\toprule
Model & CDR-H1 & CDR-H2 & CDR-H3 \\
\midrule
AntiBERTy [P]$^{\dagger}$ & \best{76.70} & \best{71.10} & 42.70 \\
AbLang2 [P]$^{\dagger}$   & 76.30 & 70.60 & 42.70 \\
Ophiuchus-Ab [P]           & 75.50 & 70.18 & \best{43.55} \\
\bottomrule
\end{tabular}
\end{table}

\paragraph{Heavy$\rightarrow$light pairing (OAS holdout500).}
Generated light chains scored by ImmunoMatch against the reference
(\Cref{tab:ab-pair}). These are local evaluations [L]; the anchor paper does
not publish p-IgGen/LICHEN values for this table. All rows use the \emph{same}
3-residue light prompt (prompt3); the prompt-free (de-novo) variants of
p-IgGen/LICHEN and \Ours{} are in \Cref{sec:ab-extra}. The p-IgGen/LICHEN
baselines are therefore not paper-transcribed. Note that \Ours{}'s mask-fill decoding reveals
the reference length, so its chain/V-gene match and diversity are not a
like-for-like comparison with the variable-length de-novo generators
(discussed in \Cref{sec:ab-extra}).

\begin{table}[t]
\centering
\caption{Antibody heavy$\rightarrow$light pairing [L] (OAS holdout500, $n{=}8$/heavy,
prompt3 for all rows). Reference ImmunoMatch ${=}0.699$; the Ophiuchus-Ab paper
reports ${\approx}0.701$ under this protocol. Best per column in \best{bold}.
$^{\ddagger}$Chain match for \Ours{} is inflated by mask-fill length revelation
(see text); not ranked against de-novo generators.}
\label{tab:ab-pair}
\small
\begin{tabular}{@{}lcccc@{}}
\toprule
Model & ImmunoMatch & gen${>}$ref & chain match & diversity \\
\midrule
p-IgGen (prompt3) & \best{0.683} & \best{0.473} & 0.990 & \best{0.368} \\
LICHEN (prompt3)  & 0.632 & 0.421 & \best{0.998} & 0.255 \\
\OursCkpt{}       & 0.600 & 0.364 & 1.000$^{\ddagger}$ & 0.118 \\
\bottomrule
\end{tabular}
\end{table}

\paragraph{FLAb property regression.}
The primary baseline is MINT Figure~3b [P], which uses nested 10-outer / 5-inner
cross-validation and reports $R^2$ (\Cref{tab:ab-flab}). The former local
5-fold pooled-Spearman sweep is a different protocol and is not a substitute.
The nested-$R^2$ runner has been repaired, but the local rerun is pending.

\begin{table}[t]
\centering
\caption{FLAb antibody property regression, paper-reported $R^2$ mean
$\pm$ standard deviation from MINT Figure~3b [P] (nested 10$\times$5-fold CV).}
\label{tab:ab-flab}
\small
\setlength{\tabcolsep}{4pt}
\begin{adjustbox}{max width=\textwidth}
\begin{tabular}{@{}lcccc@{}}
\toprule
Model [P] & trastuzumab & d44 & g6 Kd & g6 ER \\
\midrule
AbLang          & $0.293\pm0.117$ & $0.246\pm0.038$ & $0.244\pm0.034$ & $0.439\pm0.027$ \\
AntiBERTy       & $0.239\pm0.102$ & $0.217\pm0.056$ & $0.199\pm0.025$ & $0.401\pm0.032$ \\
IgBert-unpaired & $0.278\pm0.094$ & $0.181\pm0.040$ & $0.177\pm0.018$ & $0.347\pm0.023$ \\
IgBert          & $0.306\pm0.114$ & $0.131\pm0.047$ & $0.174\pm0.032$ & $0.400\pm0.023$ \\
IgT5-unpaired   & $0.299\pm0.119$ & $0.245\pm0.040$ & $0.206\pm0.029$ & $0.567\pm0.025$ \\
IgT5            & $0.274\pm0.070$ & $0.297\pm0.057$ & $0.179\pm0.014$ & $0.548\pm0.067$ \\
MINT            & $0.398\pm0.078$ & $0.379\pm0.058$ & $0.253\pm0.041$ & $0.657\pm0.023$ \\
\bottomrule
\end{tabular}
\end{adjustbox}
\end{table}

%--------------------------------------------------------------------
\subsection{Nanobody multi-task (NbBench)}
\label{sec:nbbench}

\Cref{tab:nbbench} transcribes NbBench Table~5 [P]: the paper's ESM2-150M row
and best of its eleven evaluated models for each of eleven tasks. The local
ESM2/k-mer/\Ours{} runs use a single-seed sklearn probe, whereas the paper uses
a learned MLP head and three seeds; they are retained only as [L] diagnostics
and are not mixed into this table. Paratope is especially important: the paper
reports AUROC, while the former local main column reported AUPRC. The local
hTNFa task is an additional twelfth task absent from Table~5.

\begin{table}[t]
\centering
\caption{NbBench paper-reported Table~5 baselines [P] (MLP head, three-seed
mean). Paratope uses AUROC. Local sklearn probes and the extra hTNFa task are
excluded because their protocols do not match.}
\label{tab:nbbench}
\small
\setlength{\tabcolsep}{3.5pt}
\begin{adjustbox}{max width=\textwidth}
\begin{tabular}{@{}llcc@{}}
\toprule
Task & Paper metric & ESM2-150M [P] & Best of 11 [P] \\
\midrule
VRClassification    & Accuracy & 0.9989 & 0.9989 (ESM2-150M / AntiBERTa2) \\
CDRInfilling        & BLOSUM62 recovery & 1.499 & 1.551 (ESM2-650M) \\
SARS-CoV-2          & AUROC    & 0.834 & 0.884 (AbLang-H) \\
hIL6                & AUROC    & 0.850 & 0.925 (AbLang-H) \\
Paratope            & AUROC    & 0.922 & 0.939 (AntiBERTa2-CSSP) \\
thermo-seq          & Spearman & 0.389 & 0.587 (AntiBERTa2) \\
thermo-tm           & Spearman & 0.301 & 0.593 (AntiBERTa2-CSSP) \\
polyreaction        & AUROC    & 0.833 & 0.842 (ESM2-650M) \\
nanobody\_type      & Accuracy & 0.994 & 0.999 (AntiBERTy) \\
vhh\_affinity-seq   & Spearman & 0.170 & 0.184 (IgBert) \\
vhh\_affinity-score & Spearman & 0.063 & 0.128 (AbLang-L) \\
\bottomrule
\end{tabular}
\end{adjustbox}
\end{table}

%--------------------------------------------------------------------
\subsection{Summary}
The audit establishes source-labelled baselines and a single local plug-in
interface. T1, T2, MINT, CDR infilling, FLAb, and NbBench now use paper values
as their primary references whenever the local protocol differs. T4 uses an
official-artifact sparse-13 re-score, and T3 remains explicitly local with a
separate paper audit. Consequently, this revision makes no cross-protocol claim
that \Ours{} exceeds the MINT, FLAb, CDR-infilling, or NbBench paper baselines;
those comparisons require paper-aligned reruns.
